\section{Síntesis y ampliación}

La semana 4 examinó el uso de los Coding Agents desde dos perspectivas complementarias: Mihail Eric explicó la metodología y las técnicas de gestión del Agent Manager, mientras que Boris Cherny compartió, desde la perspectiva del creador de Claude Code, la filosofía de diseño de herramientas y su experiencia práctica.

\subsection{Puntos clave}

\begin{itemize}
    \item El desarrollo de software está pasando de «personas que gestionan personas» a «personas que gestionan Agentes»: esta es la cuarta transformación de la forma organizacional
    \item Cuatro técnicas para dirigir a los Agentes: archivos de comportamiento, hooks, commands, subagents
    \item Prácticas clave: barreras de protección (pruebas + CI/CD), marcas de auditoría, correspondencia de modelos, checkpoints frecuentes
    \item Filosofía de diseño de Claude Code: terminal-native, acceso de bajo nivel, infinitamente modificable
    \item Elegir el flujo de trabajo según el tipo de tarea: explore-plan-code, TDD, iteración visual
    \item Prepararse para la capacidad de los modelos dentro de 6 meses
\end{itemize}

\subsection{Lista de implementación}

\begin{enumerate}
    \item Establecer commands, hooks y archivos de comportamiento reutilizables para las tareas comunes
    \item Usar subagents para tareas paralelas con división de trabajo clara, en lugar de amontonar todo el trabajo en un solo contexto
    \item Tratar las pruebas, la CI y la revisión humana como límites duros del flujo de trabajo del Agente, no como remedios posteriores
\end{enumerate}

\subsection{Tabla comparativa de división de roles}

\begin{table}[H]
\centering
\begin{tabular}{p{0.22\textwidth} p{0.28\textwidth} p{0.28\textwidth}}
\toprule
Rol & Responsabilidad principal & Forma de remediar el fallo \\
\midrule
Agent Manager humano & Definir objetivos, dividir tareas, establecer límites, tomar la decisión final & Redistribuir tareas, agregar restricciones, decidir si revertir o continuar \\
Agente ejecutor & Leer el contexto, modificar el código, ejecutar pruebas, generar commits & Exponer problemas mediante hooks / pruebas / revisión y reintentar \\
Sistema de revisión y CI & Ofrecer verificación objetiva, bloquear que errores evidentes lleguen a la rama principal & Reducir el riesgo oculto con verificaciones estandarizadas y bitácoras reproducibles \\
\bottomrule
\end{tabular}
\caption{Límites de responsabilidad entre humanos y Agentes en el modo Agent Manager}
\end{table}

\subsection{Modelo de madurez de adopción por equipos}

Del contenido del curso se puede extraer una ruta de adopción muy práctica:

\begin{enumerate}
    \item \textbf{Etapa de experimentación individual}: ingenieros particulares usan el Agente para responder preguntas, generar código de andamiaje y hacer refactorizaciones puntuales
    \item \textbf{Etapa de normalización de equipo}: se empiezan a consolidar commands, hooks, reglas de revisión y archivos de comportamiento unificados
    \item \textbf{Etapa de colaboración en paralelo}: el Agente se usa en tareas de múltiples líneas con división de trabajo clara, y los desarrolladores comienzan a asumir el rol de Agent Manager
    \item \textbf{Etapa de integración sistémica}: el flujo de trabajo del Agente se integra a fondo con la CI, la revisión de código, la base de conocimiento y las herramientas de gestión de proyectos
\end{enumerate}

\begin{knowledgebox}{La señal de que aumenta la madurez no es «usar más Agentes», sino «que los fallos sean más controlables»}
Los equipos que en verdad alcanzan una etapa madura no son los que se ven más vistosos, sino aquellos en los que, cuando el Agente comete un error, las pruebas, los hooks, la revisión de código y los límites de responsabilidad lo detectan a tiempo. Por eso tanto Boris como Mihail insisten repetidamente en las barreras de protección, y no solo en la velocidad.
\end{knowledgebox}

\subsection{Lecturas adicionales}

\begin{itemize}
    \item Sitio oficial de Claude Code: \url{https://claude.ai/code}
    \item Documentación del SDK de Claude Code: \url{https://docs.anthropic.com/en/docs/claude-code}
    \item Awesome Claude Agents: \url{https://github.com/vijaythecoder/awesome-claude-agents}
    \item SuperClaude Framework: \url{https://github.com/SuperClaude-Org/SuperClaude_Framework}
\end{itemize}

%% --- Fin del contenido --- %%

\end{document}
